\chapter{相关工作}
\label{chap:related}

\section{持续学习与灾难性遗忘}

灾难性遗忘\upcite{D1}是神经网络在序列化任务上训练的固有难题：在新任务上的梯度更新会覆盖对旧任务重要的参数。缓解思路大致分三类：\textbf{正则化}（约束重要参数漂移，如 EWC\upcite{A11}、梯度情景记忆 GEM\upcite{D1}）、\textbf{参数隔离}（为不同任务分配不同子网络）、以及\textbf{经验回放}（保留并复用旧数据）。本文属于回放路线，并叠加一路轻量 KL 正则作为补充锚定。

\section{经验回放及其改进}

经验回放\upcite{A2}是抗遗忘最实用、最广泛采用的手段：维护一个历史经验缓冲区，训练新数据时混入旧样本一同更新。CLEAR\upcite{A2}以均匀回放加行为克隆稳定持续强化学习，是本文的重要基线。后续工作从多个维度改进缓冲区：优先级采样\upcite{A9}、分区管理、压缩与核心集、生成式回放、自适应缓冲区。改进版蓄水池采样进一步平衡了旧经验的长期保留概率。

这些工作共享一个隐含前提：被回放的经验携带忠实的学习信号。本文的缓冲区设计在此基础上强调\textbf{按能力分区、禁止跨桶挤占}与\textbf{按能力空间距离决定回放权重}，以适配多能力域顺序训练的场景。

\section{大模型的强化微调与持续学习}

强化学习已成为大模型对齐与能力提升的主流手段。GRPO（Group Relative Policy Optimization）通过对同一查询的一组 rollout 做组内奖励归一化来估计优势，省去价值网络，是当前智能体强化学习的主力算法。近期工作从理论上揭示了 GRPO 与 DPO 的联系\upcite{B7}。在持续学习角度，已有研究发现强化微调天然比监督微调更不易遗忘\upcite{B2}。多轮智能体强化学习则暴露出策略多样性坍缩的 Echo Trap 现象\upcite{B4}，本文以固定开启的熵正则作为预防。

\section{智能体的持续学习}

面向 GUI／数字环境的智能体持续学习是新兴方向：CGL\upcite{C1}通过强化微调推进持续 GUI 学习。这些工作多集中于 GUI 模态与特定环境；本文则面向\textbf{纯文本、跨能力域}的通用智能体，以能力桶组织持续学习并给出按桶抗遗忘评测。
